%& -shell-escape
% !TeX root = thesis-text.tex
% !TeX encoding = UTF-8
% !TeX program = XeLaTeX

\documentclass{.local/lib/classes/ndvr-super-article-standalone}

\title[NDVR thesis-abstract]{%
    Проблема поиска нечетких дубликатов видео%
}

\begin{document}

    \section{Защита авторских прав}

    Количество размещенного в Интернете видео стремительно растет.
    Согласно данным Statista\footnote{
        Statista~— немецкая компания, занимающаяся
        сбором и анализом данных о развитии различных сфер мировой
        и региональных экономик. Официальный сайт: 
            \href{https://www.statista.com/}{statista.com}.
    } в мае 2017 года только пользователи YouTube загружали 
    1 миллиард часов видео каждый день \cite{Statista:2017}.
    Большая часть этих видео~— оригинальные материалы,
    снятые самими пользователями.
    Однако среди них были и незаконные «пиратские» копии лицензионных видео.
    При этом копии может отличаться от оригинала
    набором различных параметров, таких 
    как~разрешение, контраст, яркость, субтитры,
    частота кадров, последовательность сцен, или~содержать
    дополнительные изображения и т.д.
    Перебрать все варианты модификаций 
    и сравнить видео как бинарные строки невозможно.
    Заранее не известно какие изменения могут быть 
    внесены в~исходное видео.

    YouTube~— не единственный сервис в Интернете, 
    позволяющий загружать пользовательское видео.
    Кроме того, существует большое количество сайтов,
    которые изначально были созданы 
    для трансляции украденных материалов.

    Для таких сервисов как YouTube,
    проблема «пиратских» копий решается 
    с помощью нечетких дубликатов.
    Однако на практике незаконные 
    материалы чаще ищутся вручную.

    Для автоматического поиска копий требуется:
    \begin{itemize}
        \item   построить базу данных лицензионных видео;
        \item   для каждого загруженного видео вычислять 
                сходство с лицензионными.
    \end{itemize}
    Размеры такой базы данных и скорость 
    поиска по ней зависят от количества видео.
    Объем загружаемых материалов,
    делает задачу поиска копий очень затратной. 

    Потому от алгоритмов поиска нечётких дубликатов
    требуется высокая скорость и низкая вычислительная сложность. 

    Защита авторских прав~— наиболее востребованная 
    область применения поиска нечётких дубликатов видео
    \cite{Assent:2009, Cherubini:2009, Shen:2007, Law-To:2007}.

\end{document}
